\documentclass[10pt,letterpaper]{article}
\usepackage[utf8]{inputenc}
\usepackage[intlimits]{amsmath}
\usepackage{amsfonts}
\usepackage{amssymb}
\usepackage{ragged2e}
\usepackage[letterpaper, margin=0.5in]{geometry}
\usepackage{graphicx}
\usepackage{cancel}
\usepackage{mathtools}
\usepackage{tabularx}
\usepackage{arydshln}
\usepackage{tensor}
\usepackage{array}
\usepackage{xcolor}
\usepackage[boxed]{algorithm}
\usepackage[noend]{algpseudocode}
\usepackage{listings}
\usepackage{textcomp}
% \usepackage[pdf,tmpdir,singlefile]{graphviz}
\usepackage{mathrsfs}
\usepackage{bbm}
\usepackage{tikz}
\usepackage{tikz-cd}
\usepackage{enumitem}
\usepackage{arydshln}
\usepackage{relsize}
\usepackage{multicol}
\usepackage{scalerel}
\usepackage{upgreek}
\usepackage{ifthen}

\usetikzlibrary{bayesnet}
\setlist{noitemsep}

%%%%%%%%%%%%%%%%%%%%%%%%%%%%%
% Formatting commands
%%%%%%%%%%%%%%%%%%%%%%%%%%%%%
\newcommand{\n}{\hfill\break}
\newcommand{\nn}{\vspace{0.5\baselineskip}\n}
\newcommand{\up}{\vspace{-\baselineskip}}
\newcommand{\hangblock}[2]{\par\noindent\settowidth{\hangindent}{\textbf{#1: }}\textbf{#1: }\nolinebreak#2}
\newcommand{\lemma}[1]{\hangblock{Lemma}{#1}}
\newcommand{\defn}[1]{\hangblock{Defn}{#1}}
\newcommand{\thm}[1]{\hangblock{Thm}{#1}}
\newcommand{\cor}[1]{\hangblock{Cor}{#1}}
\newcommand{\prop}[1]{\hangblock{Prop}{#1}}
\newcommand{\ex}[1]{\hangblock{Ex}{#1}}
\newcommand{\exer}[1]{\hangblock{Exer}{#1}}
\newcommand{\fact}[1]{\hangblock{Fact}{#1}}
\newcommand{\remark}[1]{\hangblock{Remark}{#1}}
\newcommand{\proven}{\;$\square$\n}
\newcommand{\problem}[1]{\par\noindent{\nolinebreak#1}\n}
\newcommand{\problempart}[2]{\par\noindent\indent{}\settowidth{\hangindent}{\textbf{(#1)} \indent{}}\textbf{(#1) }\nolinebreak#2\n}
\newcommand{\ptxt}[1]{\textrm{\textnormal{#1}}}
\newcommand{\inlineeq}[1]{\centerline{$\displaystyle #1$}}
\newcommand{\pageline}{\up\par\noindent\rule{\textwidth}{0.1pt}}

%%%%%%%%%%%%%%%%%%%%%%%%%%%%%
% Math commands
%%%%%%%%%%%%%%%%%%%%%%%%%%%%%
% Set Theory
\newcommand{\card}[1]{\left|#1\right|}
\newcommand{\set}[1]{\left\{#1\right\}}
\newcommand{\setmid}{\;\middle|\;}
\newcommand{\ps}[1]{\mathcal{P}\left(#1\right)}
\newcommand{\pfinite}[1]{\mathcal{P}^{\ptxt{finite}}\left(#1\right)}
\newcommand{\naturals}{\mathbb{N}}
\newcommand{\N}{\naturals}
\newcommand{\integers}{\mathbb{Z}}
\newcommand{\Z}{\integers}
\newcommand{\rationals}{\mathbb{Q}}
\newcommand{\Q}{\rationals}
\newcommand{\reals}{\mathbb{R}}
\newcommand{\R}{\reals}
\newcommand{\complex}{\mathbb{C}}
\newcommand{\C}{\complex}
\newcommand{\halfPlane}{\mathbb{H}}
\let\HSym\H
\let\H\relax
\newcommand{\H}{\halfPlane}
\newcommand{\comp}{^{\complement}}
\DeclareMathOperator{\Hom}{Hom}
\newcommand{\Ind}{\mathbbm{1}}
\newcommand{\cut}{\setminus}
\DeclareMathOperator{\elem}{elem}

% Differentiable Manifolds
\newcommand{\RP}{\mathbb{RP}}
\newcommand{\CP}{\mathbb{RP}}
\newcommand{\osubset}{\overset{\mathclap{\scalebox{0.5}{\ptxt{open}}}}{\subset}}
\newcommand{\osubseteq}{\overset{\mathclap{\scalebox{0.5}{\ptxt{open}}}}{\subseteq}}
\newcommand{\osupset}{\overset{\mathclap{\scalebox{0.5}{\ptxt{open}}}}{\supset}}
\newcommand{\osupseteq}{\overset{\mathclap{\scalebox{0.5}{\ptxt{open}}}}{\supseteq}}
\newcommand{\pdat}[3]{\left.\pd{#1}{#2}\right|_{#3}}
\DeclareMathOperator{\so}{so}
\DeclareMathOperator{\codim}{codim}
\DeclareMathOperator{\Diff}{Diff}
\let\dSym\d
\let\d\relax
\newcommand{\d}{\partial}
\DeclareMathOperator{\gl}{gl}

% Graph Theory
\let\deg\relax
\DeclareMathOperator{\deg}{deg}
\newcommand{\degp}{\ptxt{deg}^{+}}
\newcommand{\degn}{\ptxt{deg}^{-}}
\newcommand{\precdot}{\mathrel{\ooalign{$\prec$\cr\hidewidth\hbox{$\cdot\mkern0.5mu$}\cr}}}
\newcommand{\succdot}{\mathrel{\ooalign{$\cdot\mkern0.5mu$\cr\hidewidth\hbox{$\succ$}\cr\phantom{$\succ$}}}}
\DeclareMathOperator{\cl}{cl}
\DeclareMathOperator{\affdim}{affdim}

% Probability
\newcommand{\parSymbol}{\P}
\newcommand{\Prob}{\mathbb{P}}
\renewcommand{\P}{\Prob}
\newcommand{\Avg}{\mathbb{E}}
\newcommand{\E}{\Avg}
\DeclareMathOperator{\Var}{Var}
\DeclareMathOperator{\cov}{cov}
\DeclareMathOperator{\Unif}{Unif}
\DeclareMathOperator{\Binom}{Binom}
\newcommand{\CI}{\mathrel{\text{\scalebox{1.07}{$\perp\mkern-10mu\perp$}}}}

% Standard Math
\newcommand{\inv}{^{-1}}
\newcommand{\abs}[1]{\left|#1\right|}
\newcommand{\ceil}[1]{\left\lceil{}#1\right\rceil{}}
\newcommand{\floor}[1]{\left\lfloor{}#1\right\rfloor{}}
\newcommand{\conj}[1]{\overline{#1}}
\newcommand{\of}{\circ}
\newcommand{\tri}{\triangle}
\newcommand{\inj}{\hookrightarrow}
\newcommand{\surj}{\twoheadrightarrow}
\newcommand{\ndiv}{\nmid}
\renewcommand{\epsilon}{\varepsilon}
\newcommand{\divides}{\mid}
\newcommand{\ndivides}{\nmid}
\DeclareMathOperator{\lcm}{lcm}
\DeclareMathOperator{\sgn}{sgn}
\newcommand{\map}[4]{\!\!\!\begin{array}[t]{rcl}#1 & \!\!\!\!\to & \!\!\!\!#2\\ {}#3 & \!\!\!\!\mapsto & \!\!\!\!#4\end{array}}
\newcommand{\bigsum}[2]{\smashoperator[lr]{\sum_{\scalebox{#1}{$#2$}}}}
\DeclareMathOperator{\gcf}{gcf}
\newcommand{\restr}[2]{\left.#1\right|_{#2}}

% Linear Algebra
\newcommand{\Id}{\textrm{\textnormal{Id}}}
\newcommand{\im}{\textrm{\textnormal{im}}}
\newcommand{\norm}[1]{\abs{\abs{#1}}}
\newcommand{\tpose}{^{T}\!}
\newcommand{\iprod}[1]{\left<#1\right>}
\newcommand{\giprod}{\iprod{\;\,,\;}}
\DeclareMathOperator{\tr}{tr}
\DeclareMathOperator{\trace}{tr}
\newcommand{\chgBasMat}[3]{\!\!\tensor*[_{#1}]{\left[#2\right]}{_{#3}}}
\newcommand{\vecBas}[2]{\tensor*[]{\left[#1\right]}{_{#2}}}
\DeclareMathOperator{\GL}{GL}
\DeclareMathOperator{\Mat}{Mat}
\DeclareMathOperator{\vspan}{span}
\DeclareMathOperator{\rank}{rank}
\newcommand{\V}[1]{\vec{#1}}
\DeclareMathOperator{\proj}{proj}
\DeclareMathOperator{\compProj}{comp}
\DeclareMathOperator{\row}{row}
\newcommand{\smallPMatrix}[1]{\paren{\begin{smallmatrix}#1\end{smallmatrix}}}
\newcommand{\smallBMatrix}[1]{\brack{\begin{smallmatrix}#1\end{smallmatrix}}}
\newcommand{\pmat}[1]{\begin{pmatrix}#1\end{pmatrix}}
\newcommand{\bmat}[1]{\begin{bmatrix}#1\end{bmatrix}}
\newcommand{\dual}{^{*}}
\newcommand{\pinv}{^{\dagger}}
\newcommand{\horizontalMatrixLine}{\ptxt{\rotatebox[origin=c]{-90}{$|$}}}
\DeclareMathOperator{\range}{range}
\DeclareMathOperator{\Symm}{Symm}
\DeclareMathOperator{\SU}{SU}

% Multilinear Algebra
\let\Lsym\L
\let\L\relax
\DeclareMathOperator{\L}{\mathscr{L}}
\DeclareMathOperator{\A}{\mathcal{A}}
\DeclareMathOperator{\Alt}{Alt}
\DeclareMathOperator{\Sym}{Sym}
\newcommand{\ot}{\otimes}
\newcommand{\ox}{\otimes}
\DeclareMathOperator{\asc}{asc}
\DeclareMathOperator{\asSet}{set}
\DeclareMathOperator{\sort}{sort}
\DeclareMathOperator{\ringA}{\mathring{A}}
\DeclareMathOperator{\Sh}{Sh}
\DeclareMathOperator{\Bil}{Bil}

% Topology
\newcommand{\closure}[1]{\overline{#1}}
\newcommand{\uball}{\mathcal{U}}
\DeclareMathOperator{\Int}{Int}
\DeclareMathOperator{\Ext}{Ext}
\DeclareMathOperator{\Bd}{Bd}
\DeclareMathOperator{\rInt}{rInt}
\DeclareMathOperator{\ch}{ch}
\DeclareMathOperator{\ah}{ah}
\newcommand{\LargerTau}{\mathlarger{\mathlarger{\mathlarger{\mathlarger{\tau}}}}}
\newcommand{\Tau}{\mathcal{T}}

% Analysis
\DeclareMathOperator{\Graph}{Graph}
\DeclareMathOperator{\epi}{epi}
\DeclareMathOperator{\hypo}{hypo}
\DeclareMathOperator{\supp}{supp}
\newcommand{\lint}[2]{\underset{#1}{\overset{#2}{{\color{black}\underline{{\color{white}\overline{{\color{black}\int}}\color{black}}}}}}}
\newcommand{\uint}[2]{\underset{#1}{\overset{#2}{{\color{white}\underline{{\color{black}\overline{{\color{black}\int}}\color{black}}}}}}}
\newcommand{\alignint}[2]{\underset{#1}{\overset{#2}{{\color{white}\underline{{\color{white}\overline{{\color{black}\int}}\color{black}}}}}}}
\newcommand{\extint}{\ptxt{ext}\int}
\newcommand{\extalignint}[2]{\ptxt{ext}\alignint{#1}{#2}}
\newcommand{\conv}{\ast}
\newcommand{\pd}[2]{\frac{\partial{}#1}{\partial{}#2}}
\newcommand{\del}{\nabla}
\DeclareMathOperator{\grad}{grad}
\DeclareMathOperator{\curl}{curl}
\let\div\relax
\DeclareMathOperator{\div}{div}
\DeclareMathOperator{\vol}{vol}

% Complex Analysis
\let\Re\relax
\DeclareMathOperator{\Re}{Re}
\let\Im\relax
\DeclareMathOperator{\Im}{Im}
\DeclareMathOperator{\Res}{Res}

% Abstract Algebra
\DeclareMathOperator{\ord}{ord}
\newcommand{\generated}[1]{\left<#1\right>}
\newcommand{\cycle}[1]{\smallPMatrix{#1}}
\newcommand{\id}{\ptxt{id}}
\newcommand{\iso}{\cong}
\DeclareMathOperator{\Aut}{Aut}
\DeclareMathOperator{\SL}{SL}
\DeclareMathOperator{\op}{op}
\newcommand{\isom}[4]{\!\!\!\begin{array}[t]{rcl}#1 & \!\!\!\!\overset{\sim}{\to} & \!\!\!\!#2\\ #3 & \!\!\!\!\mapsto & \!\!\!\!#4\end{array}}
\newcommand{\F}{\mathbb{F}}
\newcommand{\acts}{\;\reflectbox{\rotatebox[origin=c]{-90}{$\circlearrowright$}}\;}
\newcommand{\disjunion}{\mathrel{\text{\scalebox{1.07}{$\perp\mkern-10mu\perp$}}}}
\DeclareMathOperator{\SO}{SO}
\DeclareMathOperator{\stab}{stab}
\DeclareMathOperator{\nullity}{nullity}
\DeclareMathOperator{\Perm}{Perm}
\DeclareMathOperator{\nsubgp}{\vartriangleleft}
\DeclareMathOperator{\notnsubgp}{\ntriangleleft}
\newcommand{\presentation}[2]{\left<#1\;\middle|\;#2\right>}
\DeclareMathOperator{\Char}{char}
\DeclareMathOperator{\fchar}{char}
\DeclareMathOperator{\triv}{triv}
\DeclareMathOperator{\reg}{reg}
\DeclareMathOperator{\std}{std}
\DeclareMathOperator{\Func}{Func}
\DeclareMathOperator{\End}{End}

% Convex Optimization
\let\sectionSymbol\S
\let\S\relax
\newcommand{\S}{\mathbb{S}}
\DeclareMathOperator{\dist}{dist}
\DeclareMathOperator{\dom}{dom}
\DeclareMathOperator{\diag}{diag}
\DeclareMathOperator{\ones}{\mathbbm{1}}
\newcommand{\minimizeOver}[3]{\begin{array}{rl}\underset{#1}{\ptxt{minimize}} & #2\\ \ptxt{subject to} & #3\end{array}}
\newcommand{\maximizeOver}[3]{\begin{array}{rl}\underset{#1}{\ptxt{maximize}} & #2\\ \ptxt{subject to} & #3\end{array}}
\newcommand{\minimizationProblem}[2]{\minimizeOver{}{#1}{#2}}
\newcommand{\maximizationProblem}[2]{\maximizeOver{}{#1}{#2}}
\newcommand{\minimizeOverUnconstrained}[2]{\begin{array}{rl}\underset{#1}{\ptxt{minimize}} & #2\end{array}}
\newcommand{\maximizeOverUnconstrained}[2]{\begin{array}{rl}\underset{#1}{\ptxt{maximize}} & #2\end{array}}
\newcommand{\minimizationUnconstrained}[1]{\minimizeOverUnconstrained{}{#1}}
\newcommand{\maximizationUnconstrained}[1]{\maximizeOverUnconstrained{}{#1}}
\DeclareMathOperator{\argmin}{argmin}
\DeclareMathOperator{\argmax}{argmax}

% Proofs
\newcommand{\st}{s.t.}
\newcommand{\unique}{!}
\newcommand{\iffdef}{\overset{\ptxt{def}}{\Leftrightarrow}}
\newcommand{\eqVertical}{\rotatebox[origin=c]{90}{=}}
\newcommand{\mapsfrom}{\mathrel{\reflectbox{\ensuremath{\mapsto}}}}
\newcommand{\mapsdown}{\rotatebox[origin=c]{-90}{$\mapsto$}\mkern2mu}
\newcommand{\mapsup}{\rotatebox[origin=c]{90}{$\mapsto$}\mkern2mu}
\newcommand{\from}{\!\mathrel{\reflectbox{\ensuremath{\to}}}}
\newcommand{\labeledeq}[1]{\overset{\mathclap{\ptxt{#1}}}{=}}
\newcommand{\eqdef}{\labeledeq{def}}

% Brackets
\newcommand{\paren}[1]{\left(#1\right)}
\renewcommand{\brack}[1]{\left[#1\right]}
\renewcommand{\brace}[1]{\left\{#1\right\}}
\newcommand{\ang}[1]{\left<#1\right>}

% Algorithms
\algrenewcommand{\algorithmiccomment}[1]{\hskip 1em \texttt{// #1}}
\algrenewcommand\algorithmicrequire{\textbf{Input:}}
\algrenewcommand\algorithmicensure{\textbf{Output:}}
\newcommand{\algP}{\ptxt{\textbf{P}}}
\newcommand{\algNP}{\ptxt{\textbf{NP}}}
\newcommand{\algNPC}{\ptxt{\textbf{NP-Complete}}}
\newcommand{\algNPH}{\ptxt{\textbf{NP-Hard}}}
\newcommand{\algEXP}{\ptxt{\textbf{EXP}}}
\DeclareMathOperator{\fl}{fl}

%%%%%%%%%%%%%%%%%%%%%%%%%%%%%
% Other commands
%%%%%%%%%%%%%%%%%%%%%%%%%%%%%
\newcommand{\flag}[1]{\textbf{\textcolor{red}{#1}}}
\let\uSym\u
\let\u\relax
\newcommand{\u}[1]{\underline{#1}}
\let\bSym\b
\let\b\relax
\newcommand{\b}[1]{\textbf{#1}}
\let\iSym\i
\let\i\relax
\newcommand{\i}[1]{\textit{#1}}
\let\scSym\sc
\let\sc\relax
\newcommand{\sc}[1]{\textsc{#1}}

%%%%%%%%%%%%%%%%%%%%%%%%%%%%%%%%%%%%%%%
% Make l's curvy in math environments %
%%%%%%%%%%%%%%%%%%%%%%%%%%%%%%%%%%%%%%%
\mathcode`l="8000
\begingroup
\makeatletter
\lccode`\~=`\l
\DeclareMathSymbol{\lsb@l}{\mathalpha}{letters}{`l}
\lowercase{\gdef~{\ifnum\the\mathgroup=\m@ne \ell \else \lsb@l \fi}}%
\endgroup

%%%%%%%%%%%%%%%%%%%%%%%%%
% Fix \vdots and \ddots %
%%%%%%%%%%%%%%%%%%%%%%%%%
\usepackage{letltxmacro}
\LetLtxMacro\orgvdots\vdots
\LetLtxMacro\orgddots\ddots

\makeatletter
\DeclareRobustCommand\vdots{%
	\mathpalette\@vdots{}%
}
\newcommand*{\@vdots}[2]{%
	% #1: math style
	% #2: unused
	\sbox0{$#1\cdotp\cdotp\cdotp\m@th$}%
	\sbox2{$#1.\m@th$}%
	\vbox{%
		\dimen@=\wd0 %
		\advance\dimen@ -3\ht2 %
		\kern.5\dimen@
		% remove side bearings
		\dimen@=\wd2 %
		\advance\dimen@ -\ht2 %
		\dimen2=\wd0 %
		\advance\dimen2 -\dimen@
		\vbox to \dimen2{%
			\offinterlineskip
			\copy2 \vfill\copy2 \vfill\copy2 %
		}%
	}%
}
\DeclareRobustCommand\ddots{%
	\mathinner{%
		\mathpalette\@ddots{}%
		\mkern\thinmuskip
	}%
}
\newcommand*{\@ddots}[2]{%
	% #1: math style
	% #2: unused
	\sbox0{$#1\cdotp\cdotp\cdotp\m@th$}%
	\sbox2{$#1.\m@th$}%
	\vbox{%
		\dimen@=\wd0 %
		\advance\dimen@ -3\ht2 %
		\kern.5\dimen@
		% remove side bearings
		\dimen@=\wd2 %
		\advance\dimen@ -\ht2 %
		\dimen2=\wd0 %
		\advance\dimen2 -\dimen@
		\vbox to \dimen2{%
			\offinterlineskip
			\hbox{$#1\mathpunct{.}\m@th$}%
			\vfill
			\hbox{$#1\mathpunct{\kern\wd2}\mathpunct{.}\m@th$}%
			\vfill
			\hbox{$#1\mathpunct{\kern\wd2}\mathpunct{\kern\wd2}\mathpunct{.}\m@th$}%
		}%
	}%
}
\makeatother

\newcommand{\B}{
	\begin{tikzpicture}
	\filldraw [fill=red, draw=black] (0, 0) rectangle (0.37, 0.45);
	\draw [line width=0.5mm, white ] (0.1,0.08) -- (0.1,0.38);
	\draw[line width=0.5mm, white ] (0.1, 0.35) .. controls (0.2, 0.35) and (0.4, 0.2625) .. (0.1, 0.225);
	\draw[line width=0.5mm, white ] (0.1, 0.225) .. controls (0.2, 0.225) and (0.4, 0.1625) .. (0.1, 0.1);
	\end{tikzpicture}
}

% Allow custom symbols as arrows in tikz-cd
\tikzset{
  symbol/.style={
    draw=none,
    every to/.append style={
      edge node={node [sloped, allow upside down, auto=false]{$#1$}}}
  }
}
% Example: \arrow[r,symbol=\cong]
% https://tex.stackexchange.com/questions/394154/how-to-include-inclusion-subgroup-relationship-in-tikz-cd-diagram

\author{Professor Alejandro Uribe-Ahumada\\ \small\i{Transcribed by Thomas Cohn}}
\title{Math 635 Lecture 2}
\date{1/22/21} % Can also use \today

\begin{document}
\maketitle
\setlength\RaggedRightParindent{\parindent}
\RaggedRight

\lemma{
	Let $M$ be a $C^{\infty}$ manifold. Assume $\forall{}p\in{}M$, we have a vector space $\mathcal{E}_{p}$ (of dimension $r$). Let $\mathcal{E}=\bigsqcup_{p\in{}M}\mathcal{E}_{p}$, and $\pi:\mathcal{E}\to{}M$ the natural projection, where $\mathcal{E}_{p}\mapsto{}p\in{}M$. Assume we're given $\set{U_{\alpha}}$, a cover of $M$, plus bijections $\chi_{\alpha}$ such that $\forall\alpha$, the following diagram commutes:
	\[
		\begin{tikzcd}[ampersand replacement=\&]
			\pi\inv(U_{\alpha}) \arrow[rr,"\chi_{\alpha}"] \arrow[dr, "\pi", swap] \&[-10pt] \&[-10pt] U_{\alpha}\times\R^{r} \arrow[dl]\\
			\& U_{\alpha}
		\end{tikzcd}
	\]
	i.e. $\chi_{\alpha}(\mathcal{E}_{p})=\set{p}\times\R^{r}$.\nn
	Observe that this gives us $\forall\alpha,\beta$ \st{} $U_{\alpha}\cap{}U_{\beta}\ne\emptyset$, a map $\tau_{\alpha,\beta}:U_{\alpha}\cap{}U_{\beta}\to\GL(r,\R)$, by:
	\[
		\begin{tikzcd}[ampersand replacement=\&]
			\&[-20pt] \pi\inv(U_{\alpha}\cap{}U_{\beta}) \arrow[dl, "\chi_{\alpha}", swap] \arrow[dr, "\chi_{\beta}"] \&[-20pt]\\
			(U_{\alpha}\cap{}U_{\beta})\times\R^{r} \arrow[rr,"\chi_{\beta}\of\chi_{\alpha}\inv"] \& \& (U_{\alpha}\cap{}U_{\beta})\times\R^{r}\\[-20pt]
			(p,v) \arrow[rr,mapsto] \& \& (p,\tau_{\alpha,\beta}(p)v)\\[-20pt]
			p\in{}U_{\alpha}\cap{}U_{\beta}\\[-20pt]
			v\in\R^{r}
		\end{tikzcd}
	\]
	This mapping is a ``change of trivialization'', like a transition map for vector bundles. The matrix of $\tau_{\alpha,\beta}$ depends on $p$, but it's always a linear map.\nn
	If, $\forall\alpha,\beta$, $\tau_{\alpha,\beta}$ is $C^{\infty}$, then there is a unique topology on $\mathcal{E}$, and a unique differentiable structure, that makes $\pi:\mathcal{E}\to{}M$ a smooth vector bundle, and each $\chi_{\alpha}$ is a diffeomorphism, i.e., a smooth local trivialization.\n
}

\par\noindent
Observe: On triple intersections, $U_{\alpha}\cap{}U_{\beta}\cap{}U_{\gamma}$, $\forall{}p$, one has $\tau_{\alpha,\beta}(p)\tau_{\beta,\gamma}(p)=\tau_{\alpha,\gamma}(p)$. This is called a ``cocycle condition''.\n

\par\noindent
Now, imageine starting with a cover $\set{U_{\alpha}}$ and $C^{\infty}$ maps $\tau_{\alpha,\beta}:U_{\alpha}\cap{}U_{\beta}\to\GL(r,\R)$ satisfying the cocycle condition. Then, if we choose, $\forall{}p$, $\mathcal{E}_{p}=\R^{r}$, and $\chi_{\alpha}=I_{n}$, then we get a vector bundle $\mathcal{E}\to{}M$. So $\set{U_{\alpha};\tau_{\alpha,\beta}}$, called a \v{C}ech cocycle, is all we need to put together a vector bundle.\n

\par\noindent
Now, we apply the lemma to construct new bundles from old ones.\n

\ex{
	Given $\mathcal{E}',\mathcal{E}''\to{}M$ vector bundles, with ranks $r'$ and $r''$, respectively, define, $\forall{}p\in{}M$, $(\mathcal{E}'\oplus\mathcal{E}'')_{p}=\mathcal{E}'_{p}\oplus\mathcal{E}''_{p}$. Consider a cover $\set{U_{\alpha}}$ of $M$, and local trivializations for both $\mathcal{E}'$ and $\mathcal{E}''$. We get, $\forall\alpha$,
	\begin{align*}
		& \chi_{\alpha}':(\pi')\inv(U_{\alpha})\to{}U_{\alpha}\times\R^{r'}\\
		& \chi_{\alpha}'':(\pi'')\inv(U_{\alpha})\to{}U_{\alpha}\times\R^{r''}
	\end{align*}
	Now define $\pi:\mathcal{E}'\oplus\mathcal{E}''\to{}M$, and
	\[
		\map{\chi_{\alpha}:\pi\inv(U_{\alpha})}{U_{\alpha}\times(\R^{r'}\oplus\R^{r''})}{(p,(v',v''))}{(p,(\chi_{\alpha}'(v'),\chi_{\alpha}''(v''))\\ v'\in\mathcal{E}_{p}'\\ v''\in\mathcal{E}_{p}''}
	\]
	We have to check that the $\tau_{\alpha,\beta}$ are smooth. Note that $\tau_{\alpha,\beta}(p)$ will be a block diagonal matrix, and the two diagonal blocks vary smoothly, since $\chi_{\alpha}$ and $\chi_{\beta}$ are diffeomorphisms. Thus, $\mathcal{E}'\oplus\mathcal{E}''\to{}M$ is a natural bundle. This is called the \u{Whitney direct sum} or \u{Whitney sum}.\n
}

\newpage
\par\noindent
Observe:
\begin{itemize}[topsep=0pt, itemsep=0pt]
	\item $\rank(\mathcal{E}'\oplus\mathcal{E}'')=\rank(\mathcal{E}')+\rank(\mathcal{E}'')$
	\item $\dim(\mathcal{E}'\oplus\mathcal{E}'')=r'+r''+n$, where $n$ is the dimension of the total space.
\end{itemize}
\up\n

\par\noindent
Similarly, we can define $\mathcal{E}'\otimes\mathcal{E}''\to{}M$ with fibers $(\mathcal{E}'\otimes\mathcal{E}'')_{p}=\mathcal{E}'_{p}\otimes\mathcal{E}''_{p}$. $\rank(\mathcal{E}'\otimes\mathcal{E}'')=r'\cdot{}r''$, and $\dim(\mathcal{E}'\otimes\mathcal{E}'')=r'\cdot{}r''+n$.\n

\subsection*{Review of Tensor Products}

\defn{
	Let $V,W$ be finite-dimensional vector spaces. Their \u{tensor product} $V\otimes{}W$ is the \u{free vector space} over $V\times{}W$, modulo an equivalence relation, $\sim$. The free vector space over $V\times{}W$ is the set of all formal finite linear combinations of pairs $(v,w)\in{}V\times{}W$, and $\sim$ is defined such that
	\[
		\left\{\begin{array}{l}
			(v_{1}+v_{2},w)\sim(v_{1},w)+(v_{2},w)\\
			(\lambda{}v,w)\sim\lambda(v,w)
		\end{array}\right.
	\]
	Notation: $\forall(v,w)\in{}V\times{}W$, $v\otimes{}w=[(v,w)]$.\n
}

\par\noindent
Claim: If $(e_{1},\ldots,e_{k})$, $(f_{1},\ldots,f_{l})$ are bases of $V$, $W$, then $\set{e_{i}\otimes{}f_{j}\mid{}i\in\set{1,\ldots,k},j\in\set{1,\ldots,l}}$ is a basis of $V\otimes{}W$.\n

\cor{
	$\dim(V\otimes{}W)=\dim(V)\cdot\dim(W)$.\n
}

\par\noindent
The universal property of $V\otimes{}W$: We have a bilinear map
\[
	\map{V\times{}W}{V\otimes{}W}{(v,w)}{v\otimes{}w}
\]
It's ``universal'' in that for any bilinear map, there's a unique linear map such that the following diagram commutes:
\[
	\begin{tikzcd}
		V\times{}W \arrow[dr,"\ptxt{bilinear map}", swap] \arrow[rr, "(v{,}w)\mapsto{}v\otimes{}w"] &[-10pt] &[-10pt] V\otimes{}W \arrow[dl, "\ptxt{linear map}"]\\
		& Z
	\end{tikzcd}
\]

\par\noindent
There are other realizations of $V\otimes{}W$.
\[
	V\otimes{}W\cong\Hom(V^{*},W)\qquad\ptxt{where}\qquad{}v\otimes{}w\mapsto\paren{\begin{array}{rcl}
		V^{*}\!\!\!\! & \to\!\!\!\! & W\\
		\alpha\!\!\!\! & \mapsto\!\!\!\! & \alpha(v)w
	\end{array}}
\]
This is completely natural -- we don't need to choose a basis.\n

\par\noindent
Note: This isomorphism, as it's currently written, actually doesn't work! In an exercise, we prove that most tensors in $V\otimes{}W$ cannot be written as $v\otimes{}w$ for any $v\in{}V,w\in{}W$ (i.e. most tensors are not ``pure tensors''). Rather, they have to be written as linear combinations of pure tensors. So it would appear that this isomoprhism doesn't work. Fortunately, this map will work if we instead consider linear combinations of pure tensors.\n

\cor{
	$V^{*}\otimes{}V^{*}\cong\Hom((V^{*})^{*},V^{*})=\Hom(V,V^{*})\overset{\mathclap{\ptxt{claim}}}{\cong}\Bil(V\times{}V,\R)$, by
	\[
		\begin{array}{r>{\!\!\!\!}c>{\!\!\!\!}l}
			\Hom(V,V^{*}) & \cong & \Bil(V\times{}V,\R)\\
			(f:V\to{}V^{*}) & \mapsto & \paren{\begin{array}{r>{\!\!\!\!}c>{\!\!\!\!}l}
				V\times{}V & \to & \R\\
				(v_{1},v_{2}) & \mapsto & (f(v_{1}))(v_{2})
			\end{array}}
		\end{array}
	\]
	(Note that $f(v_{1})\in{}V^{*}$.)\nn
	Thus, $V\otimes{}V\cong\Bil(V^{*}\times{}V^{*},\R)$. This is the realization we will use!
}

\par\noindent
Note: When taking multiple tensor products:
\[
	\underbrace{V\otimes\cdots\otimes{}V}_{k}\otimes\underbrace{V^{*}\otimes\cdots\otimes{}V^{*}}_{l}\cong\{\ptxt{multilinear maps }\underbrace{V^{*}\times\cdots\times{}V^{*}}_{k}\times\underbrace{V\times\cdots\times{}V}_{l}\to\R\}
\]
Now, we put everything together, and define tensor bundles.\n

\defn{
	Given a smooth manifold $M$, $\forall{}p\in{}M$, we define
	\[
		T^{(k,l)}(T_{p}M)=\underbrace{T_{p}M\otimes\cdots\otimes{}T_{p}M}_{k}\otimes\underbrace{T_{p}^{*}M\otimes\cdots\otimes{}T_{p}^{*}M}_{l}
	\]
	$k$ is called the \u{contravariant degree}, and $l$ is called the \u{covariant degree}. Using the lemma from before, we get \u{tensor bundles}
	\[
		\bigsqcup_{p\in{}M}T^{(k,l)}(T_{p}M)=T^{(k,l)}(TM)\to{}M
	\]
}

\exer{
	Compute the rank of this bundle.\n
}

\par\noindent
In coordinate $(x^{1},\ldots,x^{n})$ on $U\osubseteq{}M$, we have a moving frame of $T^{(k,l)}(TM)$:
\[
	\set{\pd{}{x^{i_{1}}}\otimes\cdots\otimes\pd{}{x^{i_{k}}}\otimes{}dx^{j_{1}}\otimes\cdots\otimes{}dx^{j_{l}}\setmid{}i_{a},j_{b}\in\set{1,\ldots,n}}
\]

\defn{
	Any smooth section of a tensor bundle is called a \u{tensor}.\n
}

\par\noindent
Note that any tensor is a combination of these basis elements with $C^{\infty}$ functions as coefficients:
\[
	\sum{}A^{i_{1}\cdots{}i_{k}}_{j_{1}\cdots{}j_{k}}\pd{}{x^{i_{1}}}\otimes\cdots\otimes\pd{}{x^{i_{k}}}\otimes{}dx^{j_{1}}\otimes\cdots\otimes{}dx^{j_{l}}
\]

\defn{
	A \u{Riemannian metric} is a symmetric, positive-definite $(0,2)$ tensor $g_{p}:T_{p}M\times{}T_{p}M\to\R$.\n
}

\end{document}